\documentclass[11pt,a4paper]{article}
\usepackage[margin=25mm]{geometry}
\usepackage{amsmath,amssymb,amsthm,booktabs,tabularx}
\usepackage[hidelinks]{hyperref}
\usepackage{array}
\setlength{\emergencystretch}{2em}
\newtheorem{proposition}{Proposition}
\newtheorem{theorem}{Theorem}
\newtheorem{corollary}{Corollary}
\newcommand{\R}{\mathbb R}
\newcommand{\E}{\mathbb E}
\newcommand{\one}{\mathbf 1}
\newcommand{\CVaR}{\operatorname{CVaR}}
\newcommand{\col}{\operatorname{col}}
\title{From a failed certificate to a risk-aware mechanism:\\
A source audit and conditional repair of Q7 through P1}
\author{Pathfinder research note}
\date{14 September 2026}
\begin{document}
\maketitle

\begin{abstract}
We check the Q7P1 investigation directly against the cached TeX and public
version-one HTML of its source papers. Q7's strict matrix certificate is
incompatible with its implementation equalities. Its explicit payment also
fails its stated participation condition, and a deterministic example gives
an equilibrium with negative participation utility. A separate printed
reduction used for learning changes the sign of own-price curvature. These
findings do not invalidate the entire mechanism: direct price optimality
conditions recover welfare implementation under a local concavity condition.
An opponent-only transfer correction then restores equilibrium budget balance
and individual rationality without changing best responses. P1 motivates
targeting a fixed expected empirical CVaR surrogate, with inward queries that
preserve the original allocation set and its zero-allocation opt-out. We give
explicit error bounds and a monotone-operator replacement for the failed
full-message certificate. This is a source-checked mathematical proposal,
not an externally reviewed paper or a validated implementation of Q7's
original stochastic learner.
\end{abstract}

\section{Scope and source identification}

Q7 is Garjani, Shokri, and Kebriaei's mechanism-design paper
\cite{q7}; P1 is Wang, Huang, Xu, and Johansson's distributed CVaR paper
\cite{p1}. Both public records list a single version at the time of this
audit. The cached sources are \path{sources/2608.29130.tex} and
\path{sources/2609.04460.tex}. The relevant displayed expressions agree
with the public version-one HTML. The PDF endpoints could not be retrieved
through the browsing tool; this is not a claim of a byte-for-byte PDF audit.

The explicit Q7 payment is numbered (24) in the public HTML; its TeX label
is \texttt{eq40}. References to ``equation 40'' in earlier working notes
should be understood as that source label, not the rendered equation number.
The existing Q7P1 follow-up already proposed an opponent-only payment
correction and a weighted monotonicity argument. This note rederives those
ideas from the sources, sharpens the scope of the negative findings, and
adds a direct price-equilibrium argument separating the static repair from
the stronger condition needed for the proposed learner. It does not claim
that this audit independently originated every ingredient.

We write $\gamma>0$ for the payment scale, reserving $m>0$ for valuation
curvature. Q7 uses the same symbol $\alpha$ for both. Let $N\geq2$, let
$x_n\in X_n\subset\R^K_{\geq0}$, and let $p_n\in\R^K_{\geq0}$. The
local sets are convex and compact and contain zero. Prices in the static
game are not capped. The capacity constraint is $\sum_nx_n\leq c$.
We assume that the planner problem has a KKT multiplier; a suitable Slater
condition suffices. Upper bounds on prices require a separate analysis and
are not silently included in the theorems below.

\section{What is wrong in Q7, as stated}

\subsection{The matrix conditions have no common solution}

Q7's Proposition 1 requires
\begin{equation}
 \Psi+\Psi^{\mathsf T}\preceq-\nu I,\qquad \nu>0,
 \label{eq:strict}
\end{equation}
where the price-price block in row $n$, column $m$ of $\Psi$ is
$-A^n_{nm}$. Its implementation condition P2(i) requires
\begin{equation}
 \sum_m A^n_{nm}=0\qquad\text{for every }n.
 \label{eq:rowsum}
\end{equation}
These are public equations (5)(ii) and (7)(i), respectively; the cached
source locations are lines 200--216 and 240--248.

\begin{proposition}[Certificate incompatibility]
No coefficients satisfy both \eqref{eq:strict} and \eqref{eq:rowsum}.
\end{proposition}
\begin{proof}
For nonzero $v\in\R^K$, take $d=\col((0,v),\ldots,(0,v))$. Then
\begin{equation}
 d^{\mathsf T}\Psi d
 =-\sum_n v^{\mathsf T}\left(\sum_m A^n_{nm}\right)v=0.
\end{equation}
Consequently $d^{\mathsf T}(\Psi+\Psi^{\mathsf T})d=0$, contradicting
the upper bound $-\nu\|d\|^2<0$.
\end{proof}

This is an algebraic contradiction, not a failed numerical search for
coefficients. It makes the combined sufficient conditions vacuous and
invalidates the assertion that the displayed example satisfies them.
It does not prove that the mechanism is unstable. Nor does it establish
$\Psi d=0$: the allocation components can respond to a common-price shift.
The obstruction is a zero quadratic form, not necessarily a null vector
of the full operator. Replacing strict negativity by a non-strict inequality
without analysing the cross terms is not a repair.

\subsection{The explicit payment violates participation}

Write $S=\sum_m x_m$, $P_{-n}=\sum_{m\ne n}p_m$, and
$H_{-n}=\sum_{m\ne n}p_m^{\mathsf T}x_m$. Q7's explicit payment is
\begin{align}
 t_n^Q=\gamma\bigg[&\frac12\|p_n\|^2
 -\frac{p_n^{\mathsf T}P_{-n}}{N-1}
 -\frac{p_n^{\mathsf T}(S-c)}{N-1}\nonumber\\
 &+\frac{N P_{-n}^{\mathsf T}(x_n-c/N)}{(N-1)^2}
 +\frac{\sum_{m\ne n}\|p_m\|^2}{2(N-1)}
 -\frac{H_{-n}-P_{-n}^{\mathsf T}c/N}{(N-1)^2}\bigg].
 \label{eq:original-tax}
\end{align}
The TeX source is lines 410--426. Expanding its linear-price terms gives
\begin{equation}
 \sum_m a_m^n=\frac{\gamma c}{N(N-1)}.
 \label{eq:failed-participation-condition}
\end{equation}
For a positive capacity coordinate this violates Q7's condition P4(ii),
which requires that sum to be nonpositive. Thus the explicit formula does
not satisfy the advertised sufficient conditions. The following example
also shows an actual failure of its claimed participation property.

\begin{proposition}[A deterministic participation counterexample]
Set $N=3$, $K=1$, $c=\gamma=1$, $X_n=[0,1]$, and
\begin{equation}
 V_1(x)=10x-\tfrac12x^2,\qquad
 V_2(x)=V_3(x)=\tfrac1{10}x-\tfrac12x^2.
 \label{eq:counter-valuations}
\end{equation}
The profile $x=(1,0,0)$, $p=(9,9,9)$ is a Nash equilibrium for
\eqref{eq:original-tax}, but agent one's utility is $-1$.
\end{proposition}
\begin{proof}
The valuations are smooth, normalised at zero, and $1$-strongly concave.
The sets are compact and convex, contain zero, and admit strict aggregate
slack. Deterministic satisfaction is a permissible special case of Q7's
stochastic model. Each player's own negative-utility Hessian is
\begin{equation}
 \begin{pmatrix}1&-1/2\\-1/2&1\end{pmatrix},
\end{equation}
which is positive definite. At the displayed profile, the price derivatives
of negative utility are zero and its allocation derivatives are
$(0,89/10,89/10)$. These satisfy the box-constrained first-order conditions,
which are globally sufficient by convexity. Direct substitution gives
\begin{equation}
 t^Q=(21/2,-21/4,-21/4),\qquad V=(19/2,0,0),
\end{equation}
so the first utility is $19/2-21/2=-1$.
\end{proof}

The example does not satisfy P4(ii), so it is not a counterexample to a
conditional participation argument whose premises include P4(ii) and the
required equilibrium structure. It refutes the explicit formula's claimed
participation guarantee and the assertion that this formula satisfies P4.
This distinction corrects the broader informal description of the failure.

The same allocation with any common price $p_n=\lambda$,
$\lambda\in[1/10,9]$, satisfies the same sufficient first-order conditions.
Thus this instance also directly refutes uniqueness of the full message
profile. The allocation is unique; the supporting price need not be.

\subsection{The printed learning reduction is not equivalent}

In Section 4, Q7 replaces $t_n^Q$ by a reduced payment
$\widehat t_n=\gamma[s_n^{\mathsf T}Ms_n+
(\Delta\bar s_{-n}+d)^{\mathsf T}s_n]$, declaring that only terms
independent of the player's own message have been removed. But its printed
$M$ has price-price block $-I/2$ (cached lines 432--440). Therefore
\begin{equation}
 \nabla^2_{p_np_n}t_n^Q=\gamma I,
 \qquad
 \nabla^2_{p_np_n}\widehat t_n=-\gamma I.
 \label{eq:sign-error}
\end{equation}
Deleting opponent-only terms cannot change this curvature. This establishes
a printed inconsistency, not a finding about unpublished code or author
intent. In particular, repairing the static transfers does not establish
the stated stochastic-learning theorem. The replacement below uses the
derivatives of the explicit payment, not this reduced expression.

\section{A static repair that preserves incentives}

Define the opponent-only correction
\begin{equation}
 r_n(s_{-n})=\frac{\gamma}{(N-1)^2}
 \left(P_{-n}^{\mathsf T}c-NH_{-n}\right),
 \qquad t_n^R=t_n^Q-r_n(s_{-n}).
 \label{eq:rebate}
\end{equation}
The correction can have either sign. Because it is independent of $(x_n,p_n)$,
it leaves every unilateral payoff difference, best response, and Nash
equilibrium unchanged for a fixed collection of valuations. A useful exact
identity is
\begin{equation}
 p_1=\cdots=p_N=p
 \quad\Longrightarrow\quad
 t_n^R=\gamma p^{\mathsf T}(x_n-c/N),
 \label{eq:common-tax}
\end{equation}
even away from aggregate feasibility.

\begin{theorem}[Allocation implementation, balance, and participation]
Suppose each $V_n$ is continuously differentiable and $m$-strongly concave,
with $V_n(0)=0$, and
\begin{equation}
 m\geq\frac{\gamma}{(N-1)^2}.
 \label{eq:local-concavity}
\end{equation}
Under the set and KKT assumptions in Section 1, the game with payments
$t_n^R$ has a Nash equilibrium. Every equilibrium implements the unique
welfare-maximising allocation, uses a common nonnegative price, balances
the budget, and gives each agent utility at least zero. Full message
uniqueness is not asserted.
\label{thm:static}
\end{theorem}

\begin{proof}
Let $h=S-c$. At any equilibrium the price first-order conditions, coordinate
by coordinate, are
\begin{equation}
 p_n\geq0,\quad
 g_n:=Np_n-\sum_jp_j-h\geq0,\quad p_ng_n=0.
 \label{eq:price-kkt}
\end{equation}
They follow from the strictly convex own-price part of either payment.
At an agent with minimum price, $Np_n-\sum_jp_j\leq0$; hence $g_n\geq0$
forces $h\leq0$. If the maximum price is positive, its complementarity
condition yields
$0=Np_{\max}-\sum_jp_j-h$. Both terms on the right are nonnegative,
so $h=0$ and all prices equal the maximum. Otherwise all prices are zero.
Applied to each resource, this proves common pricing, feasibility, and
$p^{\mathsf T}(S-c)=0$ at every equilibrium.

The own-allocation derivative of the payment is
\begin{equation}
 \nabla_{x_n}t_n^R=
 \gamma\left(-\frac{p_n}{N-1}
             +\frac{N P_{-n}}{(N-1)^2}\right).
 \label{eq:allocation-gradient}
\end{equation}
At common prices this is $\gamma p$. The allocation conditions are
therefore exactly the planner KKT conditions with multiplier $\lambda=\gamma p$.
Strict concavity makes the planner's allocation unique.

Conversely, take a planner optimum and a KKT multiplier $\lambda$ and set
$p_n=\lambda/\gamma$. The own negative utility, up to affine and opponent-only
terms, can be written as
\begin{equation}
 -V_n(x_n)-\frac{\gamma}{2(N-1)^2}\|x_n\|^2
 +\frac{\gamma}{2}\left\|p_n-\frac{x_n}{N-1}\right\|^2.
\end{equation}
It is convex under \eqref{eq:local-concavity}. The planner KKT conditions
thus give globally optimal unilateral choices, proving equilibrium existence.

Finally, \eqref{eq:common-tax} and complementarity give
$\sum_n t_n^R=\gamma p^{\mathsf T}(S-c)=0$. An agent can deviate to
$(0,p)$ while the other messages remain fixed. This is locally feasible
and gives utility $\gamma p^{\mathsf T}c/N\geq0$. Its equilibrium utility
is no smaller. In the counterexample, the corrected taxes are $(6,-3,-3)$
and utilities are $(7/2,3,3)$.
\end{proof}

The price argument also applies to Q7's original payment. Under the same
concavity and KKT premises, its equilibrium allocations can therefore be
rescued without its impossible certificate. The transfer correction repairs
participation, not allocation incentives. Budget balance is an equilibrium
property; neither payment guarantees physically feasible requests at every
off-equilibrium message profile. The proof also explains why unique
allocation is the appropriate general conclusion, rather than unique prices.

\section{How P1 induces a different mechanism-design target}

P1's fixed-batch analysis identifies a deterministic object: the expected
empirical, smoothed CVaR, rather than the population CVaR itself. Its
one-point estimator is unbiased for the gradient of that surrogate
(Section 4.3, cached lines 853--878). This is the useful connexion:
design the strategic mechanism for the objective its oracle actually
estimates, and then control the difference from the scientific target.

There are genuine interface changes. P1 concerns cooperative agents sharing
a decision vector; Q7 has selfish agents with separate resource demands.
P1 assumes a ball about zero inside its decision set, whereas a nonnegative
resource set generally has zero on its boundary. Translating or shrinking
the strategic allocation set would change the meaning of zero consumption
and the feasible welfare benchmark. We instead move the \emph{queries},
not the allocations.

\subsection{A fixed surrogate with feasible queries}

For agent $n$, let
\begin{equation}
 C_n(x)=\CVaR_{\beta_n}[J_n(x,\xi)],\qquad 0<\beta_n<1.
\end{equation}
Assume samplewise convexity, $L_n$-Lipschitz continuity, and
$|J_n(x,\xi)|\leq U_n$ on $X_n$. These are P1-type oracle assumptions,
now applied to separate local resource variables. Assume a known interior
ball $a_n+\rho_n\mathbb B\subset X_n$. Choose
$0<\varepsilon_n<1$ and $0<\delta_n<\varepsilon_n\rho_n$ and define
\begin{equation}
 q_n(x,u)=(1-\varepsilon_n)x+\varepsilon_na_n+\delta_nu.
 \label{eq:inward}
\end{equation}
Every query with $\|u\|\leq1$ is in $X_n$ by convexity. In particular,
the message $x=0$ remains feasible, without querying outside the resource
domain. Lower-dimensional sets require a relative-interior construction;
the present statement assumes a full-dimensional ball.

Let $\widehat C_{n,s_n}$ denote empirical CVaR from a fresh i.i.d. batch of
fixed size $s_n$. Freeze $s_n$, its distribution, and the query parameters.
Define
\begin{equation}
 \widetilde C_n(x)=
 \E_{\nu,\xi^{1:s_n}}\widehat C_{n,s_n}(q_n(x,\nu))
 +\frac{\kappa_n}{2}\|x\|^2,
 \quad \nu\sim\operatorname{Unif}(\mathbb B),
 \label{eq:surrogate}
\end{equation}
and $\widetilde V_n(x)=\widetilde C_n(0)-\widetilde C_n(x)$.
The optional $\kappa_n\geq0$ supplies missing curvature. If every sample
loss is already $\mu_n$-strongly convex, the surrogate is strongly convex
with modulus at least $\mu_n(1-\varepsilon_n)^2+\kappa_n$; for merely
convex losses take $\mu_n=0$. This follows because empirical CVaR is a
maximum of nonnegative weighted sample losses whose weights sum to one.
Smoothing and expectation preserve the modulus after the affine contraction.

Select $\gamma$ so that the minimum surrogate modulus satisfies
\eqref{eq:local-concavity}. Theorem~\ref{thm:static} now applies exactly to
the game with valuations $\widetilde V_n$. No claim that P1's cooperative
network algorithm learns this strategic equilibrium is needed or made.

\subsection{What is lost relative to true risk}

Put $R_n=\sup_{x\in X_n}\|x-a_n\|$, $D_n=\sup_{x\in X_n}\|x\|$, and
\begin{equation}
 e_n=\frac{U_n}{\beta_n}\sqrt{\frac{2\pi}{s_n}}.
\end{equation}
The empirical minimum has downward bias:
$\E\widehat C_{n,s_n}(y)\leq C_n(y)$. Its expected absolute error is at
most $e_n$, by the bounded-loss empirical-CDF estimate used in P1.
Consequently, with $b_n(x)=\widetilde C_n(x)-C_n(x)$,
\begin{equation}
 -e_n-L_n\varepsilon_nR_n
 \leq b_n(x)\leq
 L_n\varepsilon_nR_n+L_n\delta_n+\frac{\kappa_nD_n^2}{2}.
 \label{eq:one-sided}
\end{equation}
For the lower bound, use Jensen's inequality for ball smoothing about
$(1-\varepsilon_n)x+\varepsilon_na_n$, then the Lipschitz displacement
bound. For the upper bound use downward empirical bias and Lipschitz
continuity. Define the interval width
\begin{equation}
 w_n=e_n+2L_n\varepsilon_nR_n+L_n\delta_n
            +\frac{\kappa_nD_n^2}{2}.
 \label{eq:width}
\end{equation}

\begin{corollary}[True-risk consequences]
Let $\widetilde s=(\widetilde x,\widetilde p)$ be a surrogate-game
equilibrium and let $x^*$ minimise $\sum_n C_n(x_n)$ on the same original
resource-feasible set. Then
\begin{align}
 \sum_n C_n(\widetilde x_n)-\sum_n C_n(x_n^*)&\leq\sum_n w_n,\\
 C_n(0)-C_n(\widetilde x_n)-t_n^R(\widetilde s)&\geq-w_n.
\end{align}
Moreover, no unilateral deviation in the true-risk game can improve agent
$n$'s utility by more than $w_n$ at $\widetilde s$.
\end{corollary}
\begin{proof}
Surrogate optimality bounds the true aggregate gap by
$\sum_n[b_n(x_n^*)-b_n(\widetilde x_n)]$, at most $\sum_n w_n$.
The difference between true and surrogate normalised valuations is
$b_n(x)-b_n(0)$, whose magnitude is at most $w_n$. For a unilateral
deviation, the normalisation constant cancels: the change in this difference
is $b_n(x_n')-b_n(\widetilde x_n)\leq w_n$. The surrogate deviation gain
is nonpositive and the tax is identical in both games.
\end{proof}

Thus equilibrium balance and implementation are exact for the declared
surrogate game; participation and incentives relative to true CVaR have
explicit errors. If agents actually optimise true CVaR rather than the
declared surrogate, $\widetilde s$ is an approximate equilibrium, not
automatically an exact one. Regularisation is also not free: its distortion
is included in $w_n$.

\section{Replacing, rather than assuming, the convergence certificate}

Let $F$ be the negative-utility pseudogradient for the surrogate game,
and group coordinates as $(x,p)$. Deriving directly from the payments gives
\begin{align}
 F_{x_n}&=\nabla\widetilde C_n(x_n)
       +\gamma\left(-\frac{p_n}{N-1}
                   +\frac{N P_{-n}}{(N-1)^2}\right),\\
 F_{p_n}&=\frac{\gamma}{N-1}
                 \left(Np_n-\sum_jp_j-S+c\right).
 \label{eq:pseudogradient}
\end{align}
Let $D=\operatorname{diag}(I,(N-1)I/N)$ and $G=DF$. Positive scaling
of the separate allocation and price blocks preserves their product-set
normal-cone conditions, so the zeros of $G+N_{\mathcal S}$ are the same
Nash equilibria. Write $\Pi=(I_N-\one\one^{\mathsf T}/N)\otimes I_K$
and $a=\gamma(2N-1)/(N-1)^2$.

\begin{proposition}[A feasible monotonicity replacement]
If the surrogate costs are $m$-strongly convex and
\begin{equation}
 m>\frac{\gamma(2N-1)^2}{4(N-1)^4},
 \label{eq:monotonicity}
\end{equation}
then $G$ is monotone. In particular, for differences $\Delta x,\Delta p$,
\begin{align}
 \langle G(s)-G(s'),s-s'\rangle
 &\geq m\|\Delta x\|^2+\gamma\|\Pi\Delta p\|^2
                 -a\langle\Pi\Delta x,\Pi\Delta p\rangle\\
 &\geq\left(m-\frac{a^2}{4\gamma}\right)\|\Delta x\|^2
       +\gamma\left\|\Pi\Delta p-
                          \frac{a}{2\gamma}\Pi\Delta x\right\|^2.
 \label{eq:monotone-bound}
\end{align}
\end{proposition}
\begin{proof}
The allocation-price block of the unscaled operator has eigenvalue
$\gamma$ on common prices and $-a$ on price disagreements. The scaled
price-allocation block is $-\gamma(\one\one^{\mathsf T}/N)\otimes I_K$,
so their common-mode contributions to the quadratic form cancel.
The scaled price-price block is $\gamma\Pi$. Strong convexity supplies
the allocation term, and completing the square gives the second inequality.
The positive allocation coefficient follows from \eqref{eq:monotonicity}.
\end{proof}

This deliberately permits a zero common-price quadratic form. It does not
contradict the obstruction, and is stronger than the local concavity
condition needed for the static theorem. A convergent replacement is the
proximal point iteration
\begin{equation}
 s^{k+1}=(I+h\mathcal A)^{-1}s^k,
 \qquad \mathcal A=G+N_{\mathcal S},\quad h>0.
 \label{eq:ppa}
\end{equation}
Here $N_{\mathcal S}$ denotes the convex normal cone. For completeness,
$\mathcal A$ is maximal monotone: subtract $m\|x_n\|^2/2$ from each
surrogate cost and include the set indicators in the resulting proper,
closed convex function. Its subdifferential is maximal monotone. The
remaining affine operator is monotone by the calculation above and is
continuous on the full space, so their sum is maximal monotone. Theorem
\ref{thm:static} supplies a zero. The resolvent is firmly nonexpansive;
its iterates are Fejer monotone with respect to the zero set, have
vanishing steps, and converge in finite dimensions to a zero
\cite{rockafellar}. Convergence to one supporting price does not imply
that all initialisations select the same price.

An inexact version retains convergence if its errors relative to the exact
resolvent are summable. This is a contract for a replacement learner, not
a validation of Q7's parallel proximal best-response iteration: the
resolvent solves a coupled variational inequality.

P1 supplies an oracle for the allocation part. With $u$ uniform on the unit
sphere in dimension $K$ and a fresh batch, set
\begin{equation}
 \widehat g_n(x)=
 (1-\varepsilon_n)\frac{K}{\delta_n}
       \widehat C_{n,s_n}(q_n(x,u))u+\kappa_nx.
 \label{eq:oracle}
\end{equation}
It is conditionally unbiased for $\nabla\widetilde C_n(x)$, and
$\|\widehat g_n(x)\|\leq
(1-\varepsilon_n)KU_n/\delta_n+\kappa_nD_n$.
This supplies a bounded stochastic oracle for inner solves. A practical
decentralised inner solver, its sample complexity, and a proof that its
errors meet the summability requirement are not supplied here. Changing
batches between calls is compatible with the fixed surrogate; changing
their size or distribution changes the target and needs another argument.

\section{Verification, novelty, and the connexion}

The accompanying \path{check_q7.py} uses rational arithmetic for the
counterexample, payment identities, gradients, price continuum examples,
and finite-dimensional instances of the monotonicity identity. It also
enumerates a two-outcome, two-sample CVaR oracle to check its expected
surrogate and inward-query gradient identity. These are reproducible
algebra checks, not measurements on Q7's transportation data or P1's sensor
network, and they do not replace the general proofs.

The source identities, counterexample, and direct equilibrium argument are
checked here without assuming the earlier internal review's verdict. This
is still a single audit, not external peer review. The only new literature
search in this audit was a targeted check for a Q7 correction and the
standard proximal-point reference; no correction was located in that bounded
search. The payment invariance principle, convex perturbation estimates,
and proximal-point convergence are established ideas. Publication novelty
of this particular synthesis has not been established.

The proposed connexion can therefore be stated precisely: \emph{P1's
sampling semantics induce a reformulation of Q7 around a fixed risk
surrogate, allocation rather than message uniqueness, and explicitly
bounded economic error.} The audit removes a false certificate and repairs
the transfer normalisation; the resulting mechanism theorem does not need
that certificate. What remains open is the efficient stochastic implementation
of the replacement learner and independent validation of the synthesis.

\appendix
\section{Audit provenance}
The SHA-256 hashes of the cached sources at the read used for this audit are:
\begin{center}\small
\begin{tabular}{l}
Q7: \texttt{2a4e5c7f74fb0350d271636307f402af3}\\
\phantom{Q7: }\texttt{b0f734341f798312ae7a6ece9bbbe76}\\[3pt]
P1: \texttt{c5631017064059edf3bb128b0bfc9177}\\
\phantom{P1: }\texttt{3601627142112e48a39dd88f6ba0d803}
\end{tabular}
\end{center}
Each hash is split across consecutive lines for typesetting; concatenate
the two hexadecimal fragments without whitespace. Existing campaign
manuscripts, review records, and status files were not changed by this audit.

\begin{thebibliography}{9}
\bibitem{q7}
S. Garjani, M. Shokri, and H. Kebriaei.
\emph{A Systematic Approach to Mechanism Design with Stochastic Dynamic Stability}.
arXiv:2608.29130v1, 2026.
\url{https://arxiv.org/html/2608.29130v1}.
\bibitem{p1}
S. Wang, K. Huang, L. Xu, and K. H. Johansson.
\emph{Distributed risk-averse optimization via CVaR}.
arXiv:2609.04460v1, 2026.
\url{https://arxiv.org/html/2609.04460v1}.
\bibitem{rockafellar}
R. T. Rockafellar.
Monotone operators and the proximal point algorithm.
\emph{SIAM Journal on Control and Optimization}, 14(5):877--898, 1976.
\url{https://doi.org/10.1137/0314056}.
\end{thebibliography}
\end{document}
